\section{Изложение}
\label{sec:izlozhenie}

\subsection{Правна уредба и правно действие на квалифицирания електронен подпис}
\label{sec:uredba}

Основният акт е eIDAS, който се прилага пряко във всички държави членки. Той различава три
вида електронни подписи. Електронен подпис са данни в електронна форма, които се добавят
към други данни в електронна форма или са логически свързани с тях и които титулярят
използва, за да се подписва (чл.~3, т.~10 eIDAS). Усъвършенстваният електронен подпис е
свързан по уникален начин с титуляря и може да го идентифицира, създаден е с данни, които
титулярят може да използва „с висока степен на доверие и единствено под свой контрол“, и
позволява да бъде открита всяка последваща промяна в подписаните данни (чл.~26 eIDAS).
Квалифицираният електронен подпис е усъвършенстван подпис, който е създаден от устройство
за създаване на квалифициран електронен подпис и се основава на квалифицирано
удостоверение (чл.~3, т.~12 eIDAS). ЗЕДЕУУ не дава собствени определения, а препраща към
регламента (чл.~13, ал.~1--3 ЗЕДЕУУ).\footnote{Закон за електронния документ и
електронните удостоверителни услуги, обн. ДВ, бр.~34 от 2001~г. До изменението, обн. ДВ,
бр.~85 от 2017~г., с което законът е съобразен с eIDAS, той се нарича Закон за
електронния документ и електронния подпис (ЗЕДЕП). За правния режим преди eIDAS вж.
Калайджиев, А., Б.~Белазелков, Г.~Димитров и др. Електронният документ и електронният
подпис~-- правен режим. София: Сиела, 2004; Димитров, Г. Право на информационните и
комуникационните технологии~-- гражданскоправни аспекти. София: Фондация „Право и
интернет“, 2014.}

Правното действие на трите вида подпис е различно. Само за квалифицирания подпис
регламентът постановява, че правната му сила „е равностойна на тази на саморъчния подпис“
(чл.~25, пар.~2 eIDAS). За останалите той установява само забрана за дискриминация:
правната им сила и допустимостта им като доказателство не могат да бъдат оспорени
единствено защото са в електронна форма или не отговарят на изискванията за квалифициран
подпис (чл.~25, пар.~1 eIDAS). Какво действие имат те, решава националното право.
Българският законодател е избрал договорния модел: правната сила на обикновения и на
усъвършенствания подпис е равностойна на тази на саморъчния, „когато това е уговорено
между страните“ (чл.~13, ал.~4 ЗЕДЕУУ). ВКС възприема същия модел и при действието на
ЗЕДЕП: законът придава значение на подписан документ на електронния документ с КЕП, но
допуска страните да се съгласят в отношенията помежду си да придадат на обикновения
електронен подпис стойността на саморъчен.\footnote{Решение №~70 от 19.02.2014~г. по
гр.~д. №~868/2012~г., ВКС, IV~г.~о., цит. по Определение №~169 от 06.04.2017~г. по ч.~т.~д.
№~672/2017~г., ВКС, I~т.~о. За проблемите при прилагането на eIDAS от българските
съдилища вж. Господинов, Д. Проблеми на електронните документи и електронните подписи при
прилагането на Регламент eIDAS от българския съд.~-- Общество и право, 2024, №~7,
с.~19--48.} Така КЕП е равен на саморъчния подпис по силата на закона и спрямо всички, а
другите два вида~-- само по силата на уговорка и само между страните по нея. От това
разграничение зависи и разпределението на доказателствената тежест
(т.~\ref{sec:tezhest}).

Писмената форма се смята за спазена, ако е съставен електронен документ, съдържащ
електронно изявление (чл.~3, ал.~2 ЗЕДЕУУ). Законът обаче не се прилага за сделките, за
които се изисква квалифицирана писмена форма, и когато държането на документа или на
екземпляр от него има правно значение (чл.~1, ал.~2 ЗЕДЕУУ). Затова КЕП не замества
нотариалната заверка там, където законът я изисква.

Особеност на българското право е разграничението между автор и титуляр на електронното
изявление. Автор е физическото лице, което в изявлението се сочи като негов извършител, а
титуляр~-- лицето, от името на което е извършено изявлението (чл.~4
ЗЕДЕУУ).\footnote{За същността на електронното изявление вж. Калайджиев, А. Относно
същността на електронното изявление и електронния документ.~-- Пазар и право, 2002, №~1.}
Терминологията не съвпада с тази на регламента. В eIDAS „титуляр на електронен подпис“ е
физическото лице, което създава подписа (чл.~3, т.~9 eIDAS), а юридическото лице не
подписва, а създава електронен печат (чл.~3, т.~24 и~25 eIDAS). Разликата има значение и
за доказването. Квалифицираният електронен печат се ползва от изрична презумпция за цялост
на данните и за точност на техния произход (чл.~35, пар.~2 eIDAS). За квалифицирания
подпис регламентът не установява такава презумпция, а само равностойност със саморъчния
подпис. Какво следва от тази равностойност в процеса, определя националното процесуално
право, в случая ГПК.

Когато физическо лице подписва като представител на юридическо лице, в удостоверението му
може да бъде вписано и юридическото лице, защото квалифицираните удостоверения могат да
съдържат допълнителни незадължителни данни (чл.~28, пар.~3 eIDAS). В практиката такива
удостоверения се наричат „професионални“. Подписва обаче физическото лице. Дали то е
могло да задължи юридическото лице, се решава по общите правила за представителството, а
не от съдържанието на удостоверението (т.~\ref{sec:proverka}).

\subsection{Механизъм на квалифицирания електронен подпис: какво доказва и какво не}
\label{sec:mehanizam}

КЕП се създава чрез асиметрична криптография. Титулярят разполага с двойка свързани
ключове. Частният ключ представлява „данните за създаване на електронен подпис“ (чл.~3,
т.~13 eIDAS) и трябва да е известен само на титуляря. Публичният ключ се включва в
удостоверението, с което доставчикът свързва данните за валидиране на подписа с
определено физическо лице (чл.~3, т.~14 и~15 eIDAS). При подписване софтуерът изчислява
хеш стойността на документа~-- кратка контролна сума, която се променя при всяка промяна
на документа, и я подписва с частния ключ. Подписът и удостоверението се вграждат в
документа, например в PDF файла, или се прилагат към него отделно \figref{fig:podpis}.

Устройството за създаване на КЕП трябва да гарантира, че данните за създаване на подписа
на практика се срещат само веднъж, не могат да бъдат извлечени и „могат да бъдат надеждно
защитени от законния титуляр на електронния подпис срещу използване от други лица“
(приложение~II, т.~1 eIDAS). Традиционно това е смарт карта или USB токен, който се
отключва с ПИН код. Все по-често обаче частният ключ се пази в сървърно устройство на
доставчика, а титулярят само дава команда за подписване, например от мобилно приложение.
С Регламент (ЕС) 2024/1183 този модел е уреден като отделна квалифицирана услуга по
управление на устройства за създаване на КЕП от разстояние. При нея доставчикът
„генерира или управлява данни за създаване на електронен подпис от името на титуляря“
(чл.~29а, пар.~1, б.~„а“ eIDAS).

\begin{figure}[H]
\centering
\begin{tikzpicture}[node distance=0.5cm]
\node[box] (doc) {Документ};
\node[box, right=of doc] (hash) {Хеш стойност\\на документа};
\node[box, right=of hash] (sign) {Подписване с частния\\ключ в устройството\\за създаване на КЕП};
\node[box, right=of sign] (out) {Подписан документ:\\подпис и квалифицирано\\удостоверение};
\draw[arr] (doc) -- (hash);
\draw[arr] (hash) -- (sign);
\draw[arr] (sign) -- (out);
\node[done, below=1.2cm of out] (integ) {Цялост: хешът на\\получения документ\\съвпада с подписания};
\node[done, left=of integ] (cert) {Удостоверението е\\квалифицирано и валидно\\към момента на подписването};
\node[done, left=of cert] (res) {Протокол от\\валидирането};
\node[court, left=of res] (not) {Не доказва кой\\е дал командата\\за подписване};
\draw[arr] (out) -- (integ);
\draw[arr] (integ) -- (cert);
\draw[arr] (cert) -- (res);
\draw[arr, dashed] (res) -- (not);
\node[font=\small\itshape, anchor=south west, inner sep=1pt] at ([yshift=0.12cm]doc.north west) {Създаване};
\node[font=\small\itshape, anchor=north west, inner sep=1pt] at ([yshift=-0.12cm]not.south west) {Валидиране (чл.~32 eIDAS)};
\end{tikzpicture}
\caption{Създаване и валидиране на квалифициран електронен подпис}
\label{fig:podpis}
\end{figure}

Проверката на подписа се нарича валидиране. При него се потвърждава, че удостоверението
към момента на подписването е било квалифицирано и валидно и е издадено от квалифициран
доставчик, че подписът е създаден от устройство за създаване на КЕП, че целостта на
подписаните данни не е застрашена и че изискванията на чл.~26 са били изпълнени към
момента на подписването (чл.~32, пар.~1 eIDAS). Квалифицираните доставчици и услугите им
се проверяват по доверителния списък, който всяка държава членка поддържа и публикува
(чл.~22 eIDAS). Положителният резултат установява три факта: документът не е променян
след подписването; подписът е създаден с частния ключ, който съответства на
удостоверението; удостоверението е било валидно в момента на подписването.

Валидирането обаче не установява кой човек е дал командата за подписване. То доказва, че е
използван определен ключ, а не че титулярят е изразил воля. Връзката между двете е правно
предположение, основано на изискването титулярят да може да използва данните за създаване
на подписа „единствено под свой контрол“ (чл.~26, б.~„в“ eIDAS). Когато предположението не
отговаря на действителността~-- удостоверението е издадено на името на друг човек, ПИН
кодът е споделен или профилът в мобилното приложение е превзет, валидирането пак ще бъде
успешно. Затова в спора за авторството техническата проверка е необходима, но не е
достатъчна.

Решаващ е и моментът на подписването. Освен ако е уговорено друго, доставчикът спира
действието на удостоверението за срок до 48 часа по искане на титуляря, без да е длъжен да
се увери в самоличността му (чл.~26, ал.~2, т.~1 ЗЕДЕУУ). Той прекратява действието на
удостоверението по искане на титуляря или когато установи, че то е издадено въз основа на
неверни данни (чл.~27, ал.~2 и~3 ЗЕДЕУУ). Регламентът нарича прекратяването „отмяна“.
Отмененото удостоверение губи валидността си от момента на отмяната и статутът му не може
да бъде възстановен (чл.~28, пар.~4 eIDAS). Отмяната се публикува в срок до 24 часа от
искането и става валидна веднага след публикуването (чл.~24, пар.~3 eIDAS). Следователно
подписът, създаден преди отмяната, остава валиден, а създаденият след нея~-- не. Часът,
който компютърът на подписващия вписва в подписа, може да бъде променен. Надеждно
доказателство за момента дава квалифицираният електронен времеви печат, който се ползва
от презумпция за точност на датата и часа и за цялост на данните (чл.~41, пар.~2 eIDAS).

\subsection{Документът с квалифициран електронен подпис като доказателство в гражданския процес}
\label{sec:dokazatelstvo}

Електронен документ е „всяко съдържание, съхранявано в електронна форма“ (чл.~3, т.~35
eIDAS; чл.~3, ал.~1 ЗЕДЕУУ). Правната му сила и допустимостта му като доказателство не
могат да бъдат оспорени единствено защото е в електронна форма (чл.~46 eIDAS). ГПК урежда
електронния документ в раздела за писмените доказателства. ВКС приема, че и електронното
писмо, което съдържа изявление с гражданскоправно значение, е електронен документ, като
„отделен е въпросът дали този документ може [да] бъде определен като подписан или
не“.\footnote{Решение №~303 от 07.10.2013~г. по гр.~д. №~3715/2013~г., ВКС, IV~г.~о.
(постановено по чл.~290 ГПК при действието на ЗЕДЕП). За електронните доказателства вж. и
Младенов, Д. Електронни доказателства.~-- De Jure, 2022, т.~13, №~1, с.~63--74.}

Електронният документ може да бъде представен възпроизведен на хартия като препис, заверен
от страната, но при поискване страната е длъжна да го представи на електронен носител
(чл.~184, ал.~1 ГПК). От 30.06.2021~г. завереният препис може да бъде представен и като
електронен образ, заверен от страната с КЕП (чл.~183, ал.~2 ГПК).\footnote{Разпоредбата е
създадена с изменението на ГПК, обн. ДВ, бр.~110 от 2020~г.} ВКС подчертава, че
представеният на хартия документ има характера на препис на частен електронен документ и
че той „визуализира само част от електронния документ“, без информацията от метаданните,
която е достъпна само когато документът е представен на електронен
носител.\footnote{Решение №~176 от 02.07.2025~г. по т.~д. №~123/2023~г., ВКС, II~т.~о.
(по чл.~290 ГПК).} За документа с КЕП това има особено значение. Подписът съществува само
в електронния файл. Разпечатката показва съдържанието и евентуално изображение на
подписа, но не позволява подписът да бъде валидиран. Ако насрещната страна не поиска
електронния носител, преписът е годно доказателство.\footnote{Решение №~164 от
16.05.2022~г. по в.~гр.~д. №~173/2022~г., Окръжен съд~-- Пазарджик.} Поиска ли го обаче,
документът трябва да бъде представен в електронна форма. Според автора, ако страната не
стори това, по аналогия с чл.~183, ал.~1 ГПК представеният препис трябва да бъде изключен
от доказателствата.

Съдът не винаги може сам да провери подписа. Ако той не разполага с технически средства и
специалисти, които да позволят възпроизвеждането на документа и „надлежната проверка на
електронния подпис“ в съдебната зала, електронни копия се предоставят на страните, а
истинността на документа може да бъде оспорена в следващото заседание (чл.~184, ал.~2
ГПК). Нормата е логична: страната не може да реши дали да оспори документа, преди да е
могла да провери подписа.

Частните документи, подписани от лицата, които са ги издали, съставляват доказателство, че
изявленията в тях са направени от тези лица (чл.~180 ГПК). Това е формалната
доказателствена сила на документа. Тя се отнася до авторството на изявлението, а не до
верността на съдържанието му.\footnote{Вж. Сталев, Ж., А.~Мингова, В.~Попова,
О.~Стамболиев, Р.~Иванова. Българско гражданско процесуално право. 10~изд. София: Сиела,
2020.} Поради равностойността по чл.~25, пар.~2 eIDAS документът с КЕП има формалната
доказателствена сила на подписан частен документ. Документ с обикновен или усъвършенстван
подпис я има само при уговорка по чл.~13, ал.~4 ЗЕДЕУУ. Иначе съдът го преценява по
вътрешно убеждение заедно с останалите доказателства (чл.~12 ГПК). Такъв документ обаче
не е „правно нищо“: ВКС приема, че ако авторството му бъде оспорено, то се установява с
други доказателствени средства, а ако не бъде оспорено в срок, не е нужно страната да
ангажира технически експертизи.\footnote{Определение №~169 от 06.04.2017~г. по ч.~т.~д.
№~672/2017~г., ВКС, I~т.~о.}

ГПК съдържа и две по-особени правила, свързани с подписания електронен документ. Първо,
съдът преценява с оглед на всички обстоятелства доказателствената сила на документ, в който
има зачерквания, изтривания, добавки и други външни недостатъци, но това правило „не се
прилага за подписания електронен документ“ (чл.~178, ал.~2 ГПК). Причината е техническа:
всяка промяна след подписването прави подписа невалиден. Подписаният електронен документ
или е непроменен, или валидирането показва, че е променен; средно положение няма. Второ,
частният документ има достоверна дата за трети лица и от деня, в който настъпи факт,
„установяващ по безсъмнен начин предхождащото го съставяне на документа“ (чл.~181, ал.~1
ГПК). Може да се приеме, че квалифицираният времеви печат е такъв факт, защото се ползва
от презумпция за точност на датата и часа (чл.~41, пар.~2 eIDAS).

\subsection{Оспорване на документ с квалифициран електронен подпис}
\label{sec:osporvane}

\subsubsection{Ред и предмет на оспорването}
\label{sec:red}

Заинтересованата страна може да оспори истинността на документа най-късно с отговора на
съдопроизводственото действие, с което той е представен, а ако документът е представен в
съдебно заседание~-- до края на заседанието (чл.~193, ал.~1 ГПК). При електронния документ
срокът се удължава до следващото заседание, ако проверката на подписа не е могла да бъде
извършена в залата (чл.~184, ал.~2 ГПК). Проверка се открива само ако страната, която е
представила документа, заяви, че желае да се ползва от него (чл.~193, ал.~2 ГПК).

Предмет на оспорването по чл.~193 ГПК е истинността на документа, т.е. дали той изхожда от
лицето, посочено като негов автор. Наказателният закон определя неистинския документ като
документ, „на който е придаден вид, че представлява конкретно писмено изявление на друго
лице, а не на това, което действително го е съставило“ (чл.~93, т.~6 НК). От истинността
трябва да се разграничава верността на съдържанието, която се преценява с другите
доказателства.

При документа с КЕП оспорването може да има три основания. Първо, подписът не е валиден,
например защото документът е променен или удостоверението е било отменено към момента на
подписването. Второ, удостоверението е издадено на името на оспорващия, но не по негово
искане, например въз основа на чужда или подправена лична карта. Трето, удостоверението е
издадено правилно, но подписът е създаден от друго лице без знанието на титуляря. Първото
основание се установява технически, а второто и третото~-- с всички доказателствени
средства. Само при първото валидирането дава отговор. При другите две то е успешно и
спорът се води за факти извън него.

\subsubsection{Разпределение на доказателствената тежест}
\label{sec:tezhest}

Общото правило е, че всяка страна трябва да установи фактите, на които основава исканията
или възраженията си (чл.~154, ал.~1 ГПК). Чл.~193, ал.~3 ГПК съдържа специално правило за
оспорването на документи. Тежестта за доказване на неистинността пада върху страната,
която оспорва документа. Когато обаче се оспорва частен документ, „който не носи подписа
на страната, която го оспорва“, тежестта за доказване на истинността пада върху страната,
която го е представила.

Решаващият въпрос е дали документът с електронен подпис „носи подписа“ на оспорващия. За
КЕП отговорът следва от чл.~25, пар.~2 eIDAS: документът, подписан с валиден КЕП, основан
на удостоверение, издадено на оспорващия, носи неговия подпис. Затова тежестта да докаже,
че не той е подписал, е на оспорващия. Обикновеният електронен подпис без уговорка по
чл.~13, ал.~4 ЗЕДЕУУ не е равностоен на саморъчния. Документът с такъв подпис не носи
подписа на страната в смисъла на чл.~193, ал.~3 ГПК и истинността му трябва да докаже
страната, която го е представила. При уговорка по чл.~13, ал.~4 ЗЕДЕУУ тежестта отново е
на оспорващия, но съществуването на уговорката трябва да докаже страната, която се
позовава на нея. Разпределението е обобщено в табл.~\ref{tab:tezhest}.

\begin{table}[H]
\caption{Правно действие и доказателствена тежест според вида на подписа}
\label{tab:tezhest}
\small
\begin{tabularx}{\textwidth}{@{}L{4.3cm}YY@{}}
\toprule
\textbf{Вид подпис} & \textbf{Правно действие} & \textbf{Кой доказва при оспорване (чл.~193, ал.~3 ГПК)}\\
\midrule
Обикновен или усъвършенстван подпис без уговорка & Допустимо доказателство, което съдът
преценява по вътрешно убеждение (чл.~25, пар.~1 eIDAS; чл.~12 ГПК) & Страната, която е
представила документа, доказва истинността му\\
\addlinespace
Обикновен или усъвършенстван подпис с уговорка & Равностоен на саморъчния подпис между
страните по уговорката (чл.~13, ал.~4 ЗЕДЕУУ) & Оспорващият доказва неистинността, след като
другата страна докаже уговорката\\
\addlinespace
Квалифициран електронен подпис & Равностоен на саморъчния подпис по силата на закона
(чл.~25, пар.~2 eIDAS) & Оспорващият доказва неистинността, след като се установи, че
подписът е валиден КЕП, основан на негово удостоверение\\
\bottomrule
\end{tabularx}
\end{table}

Разпределението има две стъпки. Страната, която представя документа, трябва първо да
установи, че той е подписан с валиден КЕП на насрещната страна. Това обикновено става с
протокол от валидиране или чрез вещо лице. СЕС приема, че вписването в удостоверението,
че подписът е квалифициран, не е достатъчно. Когато квалификацията се оспорва, националният
съд е длъжен да провери дали са изпълнени всички кумулативни условия на чл.~3, т.~12
eIDAS, включително изискванията на чл.~26 и на приложение~I. Ако липсва квалифицирано
удостоверение, подписът не е КЕП, дори да е обозначен като „професионален електронен
подпис“.\footnote{Решение на СЕС от 20 октомври
2022~г. по дело C-362/21, \textit{Екофрукт}, ECLI:EU:C:2022:815, т.~2 и~3 от
диспозитива. Запитването е отправено от Административен съд~-- Велико Търново.} Едва след
като се установи валиден КЕП, тежестта по чл.~193, ал.~3, изр.~1 ГПК преминава върху
оспорващия. Ако валидирането е отрицателно, документът не носи КЕП и се прилага
изречение второ.

Тежестта на оспорващия обаче е трудна за изпълнение. Той трябва да докаже отрицателен
факт~-- че не е подписал, а доказателствата обикновено не са у него, а у доставчика. Това
са искането за издаване на удостоверението, копието на документа за самоличност, записът от
дистанционната идентификация и данните за командите за подписване. Регламентът задължава
квалифицирания доставчик да записва и съхранява цялата имаща отношение информация за
издадените и получените данни, включително „с оглед предоставяне на доказателство при
съдебни производства“ (чл.~24, пар.~2, б.~„з“ eIDAS). ГПК дава и процесуалното средство:
страната може да поиска от съда да задължи неучастващо в делото лице да представи
намиращ се у него документ, а третото лице, което неоснователно не го представи, отговаря
пред страната за вредите (чл.~192, ал.~1 и~3 ГПК). Ако документът е у насрещната страна,
се прилага чл.~190 ГПК, а съдът може да приеме за доказани фактите, за чието установяване
страната е създала пречки (чл.~161 ГПК).

В публикуваната практика на ВКС не беше открито решение по оспорване на истинността на
документ с КЕП. Цитираните актове се отнасят до документи без КЕП, при които въпросът за
тежестта се решава по изречение второ на чл.~193, ал.~3 ГПК.

\subsubsection{Ограниченията по чл.~18а ЗЕДЕУУ}
\label{sec:18a}

Чл.~18а ЗЕДЕУУ, озаглавен „Оспорване“, стеснява възможността за оспорване още повече.
Лицето, посочено като титуляр на електронен подпис, не може да оспори авторството и
титулярството спрямо адресата, когато електронният документ е отправен чрез информационна
система, предназначена да работи в автоматичен режим, или е създаден от лице, на което е
предоставен достъп до начина на идентифициране (ал.~1). Второто ограничение отпада от
деня, в който адресатът получи уведомление, че документът не изхожда от автора или от
титуляря (ал.~2). И двете не се прилагат, когато адресатът не е положил дължимата грижа
(ал.~3). Публикувана съдебна практика по прилагането на разпоредбата не беше открита.

Според автора разпоредбата не е правило за доказателствената тежест, а материалноправно
правило за разпределение на риска. В посочените случаи изявлението се счита за изявление
на титуляря, дори ако физически го е направил друг. Нормата следва логиката на чл.~7
ЗЕДЕУУ, според който титулярят носи риска от грешки при предаването на изявлението, освен
ако адресатът не е положил дължимата грижа. От текста следват три извода.

Първо, ограничението се отнася до всички видове електронни подписи, а не само до КЕП, и
защитава само адресата, не и трети лица.

Второ, „предоставен достъп до начина на идентифициране“ означава, че титулярят доброволно
е дал на друго лице картата и ПИН кода си или достъпа до профила си в мобилното
приложение, например на счетоводител или служител. Кражбата на устройството или
превземането на профила без знанието на титуляря не попадат тук. В тези случаи титулярят
може да оспори авторството, но тежестта е негова по чл.~193, ал.~3 ГПК.

Трето, уведомлението по ал.~2 трябва да бъде отправено до адресата, а не до доставчика.
Затова титулярят, който научи, че подписът му се използва от другиго, трябва едновременно
да поиска спиране или прекратяване на удостоверението (чл.~26, ал.~2, т.~1 и чл.~27,
ал.~2 ЗЕДЕУУ) и да уведоми известните му адресати.

\subsubsection{Проверка на истинността}
\label{sec:proverka-istinnost}

Съдът проверява истинността чрез сравняване с други безспорни документи, чрез разпит на
свидетели или чрез вещи лица (чл.~194, ал.~1 ГПК). Сравняването на почерка, което е
основният способ при саморъчния подпис, при КЕП е безпредметно. Затова основното
доказателствено средство е вещото лице, което се назначава, когато за изясняване на
спорните въпроси са необходими специални знания (чл.~195, ал.~1 ГПК). Компютърно-техническата
експертиза може да установи:
\begin{itemize}
\item дали подписът е валиден и документът не е променян след подписването;
\item дали удостоверението е издадено от доставчик, вписан в доверителния списък, и какъв
  е бил статутът му към момента на подписването;
\item дали има квалифициран времеви печат и какъв момент удостоверява;
\item въз основа на сведенията от доставчика~-- как е проверена самоличността при
  издаването на удостоверението и откъде и с какво устройство е дадена командата за
  подписване.
\end{itemize}

Последното е особено важно при дистанционното подписване. Тогава данните за създаване на
подписа са у доставчика и само той разполага с журнала на използването им. Свидетелските
показания имат допълнително значение, например за това дали титулярят е предоставял
картата или ПИН кода си на друго лице (чл.~18а, ал.~1, т.~2 ЗЕДЕУУ).

Ако съдът приеме, че документът е неистински, той го изключва от доказателствата и го
изпраща на прокурора заедно с определението си (чл.~194, ал.~2 ГПК). Съдът може да се
произнесе по оспорването и с решението по делото (чл.~194, ал.~3 ГПК).

\subsection{Злоупотреби с квалифицирания електронен подпис и разпределение на риска}
\label{sec:zloupotrebi}

\subsubsection{Удостоверение, издадено без надлежна проверка}
\label{sec:proverka}

Злоупотребата често започва още при издаването на удостоверението. Затова при издаването
на квалифицирано удостоверение доставчикът проверява самоличността и, ако е приложимо,
специфичните атрибути на лицето (чл.~24, пар.~1 eIDAS). С Регламент (ЕС) 2024/1183 методите
за проверка на самоличността са изброени изчерпателно: европейският портфейл за цифрова
самоличност или нотифицирано средство за електронна идентификация с ниво на осигуреност
„високо“; удостоверение за КЕП или за квалифициран електронен печат; други методи, които
гарантират идентифицирането „с висока степен на надеждност“ и чието съответствие се
потвърждава от орган за оценяване на съответствието; физическото присъствие на лицето
(чл.~24, пар.~1а eIDAS). Сходно е уредена проверката на атрибутите, към които може да се
отнесе и представителната власт (чл.~24, пар.~1б eIDAS). Референтните стандарти са определени с
Регламент за изпълнение (ЕС) 2025/1566, който ще се прилага от 19~август
2027~г.\footnote{Регламент за изпълнение (ЕС) 2025/1566 на Комисията от 29 юли 2025~г. за
определяне на правила за прилагането на Регламент (ЕС) №~910/2014 на Европейския
парламент и на Съвета по отношение на референтните стандарти за проверка на самоличността
и атрибутите на лице, на което ще бъде издадено квалифицирано удостоверение или
квалифицирано електронно удостоверение за атрибути, ОВ L, 2025/1566, 30.07.2025~г.,
чл.~2. Приложението към него препраща към стандарта ETSI TS 119 461 V2.1.1 (2025-02) с
някои адаптации. За
електронната идентификация вж. и Димитров, Г., С.~Ракшиева. Електронна идентификация и
електронни удостоверителни услуги.~-- Юридически сборник, 2016, т.~23, с.~100--112.}
ЗЕДЕУУ препраща към чл.~24, пар.~1 eIDAS (чл.~25, ал.~1 ЗЕДЕУУ). Нарушението на
задълженията по чл.~24, пар.~1--4 eIDAS се наказва с глоба, а за юридическите лица~-- с
имуществена санкция от 5000 до 100\,000~лв. (чл.~45, ал.~1, т.~2 и ал.~2
ЗЕДЕУУ).\footnote{Сумите са посочени така, както са в текста на закона.}

Значението на проверката личи от казуса, разгледан в курса. През 2013~г.
„Борика-Банксервиз“~АД издава професионално удостоверение за КЕП с автор управителя
П.~С.~Я. и титуляр „Автогара-Хасково“~ЕООД. Искането за издаване е подписано само от
П.~С.~Я. Според вписването в търговския регистър обаче дружеството се представлява само
заедно от двама от тримата управители, единият от които задължително е А.~Х.~А. След
жалба на последния председателят на КРС налага на доставчика имуществена санкция от
5000~лв. за нарушение на чл.~25, ал.~3, т.~1 ЗЕДЕП.\footnote{Наказателно постановление
№~123 от 07.07.2014~г. на председателя на КРС.} Районният съд потвърждава наказателното
постановление, а административният съд оставя решението му в сила.\footnote{Решение
№~302 от 22.10.2014~г. по а.н.д. №~802/2014~г., Районен съд~-- Хасково; Решение №~22 от
10.02.2015~г. по к.а.н.д. №~374/2014~г., Административен съд~-- Хасково.} Касационният съд
отхвърля довода, че волята на дружеството е изразена с договора за удостоверителни услуги,
подписан от двама управители. Според съда предмет на проверката не е договорът, а
предхождащото го искане за издаване на удостоверението. За него законът изисква то да
изхожда „от титуляря или надлежно овластено от него лице“, а управителят П.~С.~Я. не може
самостоятелно да представлява дружеството.

Изводите от казуса са валидни и днес. Чл.~25, ал.~3 ЗЕДЕП е отменен през 2017~г., но
същото задължение следва от чл.~24, пар.~1 и~1б eIDAS и от чл.~25, ал.~1 ЗЕДЕУУ, според
който квалифицираното удостоверение се издава по искане на титуляря. Казусът показва и
друго: удостоверението отразява представителната власт, проверена от доставчика към
момента на издаването, но не я създава. Ако подписалият няма представителна власт,
договорът не обвързва юридическото лице, освен ако то го потвърди във формата, предвидена
за упълномощаването, а действалият без власт дължи обезщетение на добросъвестната
насрещна страна (чл.~42 ЗЗД). Чл.~25а ЗЕДЕУУ добавя още едно правило. Титулярят трябва
незабавно да уведоми доставчика за промени в обстоятелствата, посочени в удостоверението,
а промените не могат да бъдат противопоставени на трети добросъвестни лица. Как това
правило се съчетава с публичността на търговския регистър, е въпрос, който излиза извън
обхвата на работата.

\subsubsection{Използване на подписа от друго лице}
\label{sec:drugo-lice}

Освен титуляря никой няма право на достъп до данните за създаване на електронния подпис
(чл.~18, т.~1 ЗЕДЕУУ). На практика обаче картата и ПИН кодът често се дават на
счетоводител, служител или член на семейството. Тогава се прилага чл.~18а, ал.~1, т.~2
ЗЕДЕУУ и титулярят не може да оспори авторството спрямо адресата, докато не го уведоми
(т.~\ref{sec:18a}). Рискът от доброволното предоставяне на подписа е изцяло негов.

Различно е положението, когато подписът е използван без знанието на титуляря: картата е
открадната, ПИН кодът е узнат или профилът в мобилното приложение е превзет. Тук
титулярят може да оспори авторството, но трябва да докаже твърдението си
(т.~\ref{sec:tezhest}). Затова при съмнение за компрометиране той трябва незабавно да
поиска спиране на удостоверението, което доставчикът извършва без проверка на
самоличността, и прекратяването му (чл.~26, ал.~2, т.~1 и чл.~27, ал.~2 ЗЕДЕУУ). Спиране
може да поиска и лице, за което според обстоятелствата е видно, че може да знае за
нарушения на сигурността на частния ключ, например представител, съдружник, служител или
член на семейството (чл.~26, ал.~2, т.~2 ЗЕДЕУУ). Подписите, създадени след публикуването
на отмяната, не са валидни (чл.~24, пар.~3 и чл.~28, пар.~4 eIDAS). Подписите, създадени
преди нея, остават валидни и могат да бъдат оспорени само по общия ред.

\subsubsection{Отговорност на доставчика на удостоверителни услуги}
\label{sec:otgovornost}

Доставчиците на удостоверителни услуги отговарят за щетите, причинени умишлено или поради
небрежност на физическо или юридическо лице поради неизпълнение на задълженията им по
регламента (чл.~13, пар.~1 eIDAS; чл.~29 ЗЕДЕУУ). С Регламент (ЕС) 2024/1183 изрично е
посочено, че обезщетение може да се търси както за имуществени, така и за неимуществени
щети. Тежестта на доказване е разпределена различно. Умисълът или небрежността на
неквалифицирания доставчик трябва да докаже ищецът, а при квалифицирания доставчик те се
предполагат, докато доставчикът не докаже, че щетите са настъпили без негова вина.
Доставчикът не отговаря за щети от използване на услугата извън ограниченията, за които е
информирал клиентите си предварително и които са разпознаваеми за трети страни (чл.~13,
пар.~2 eIDAS). За да покрие риска, квалифицираният доставчик поддържа достатъчни финансови
ресурси или застраховка за отговорност (чл.~24, пар.~2, б.~„в“ eIDAS). Титулярят от своя
страна отговаря спрямо доставчика, ако е предоставил неверни данни или е премълчал данни,
имащи отношение към съдържанието или към издаването на удостоверението (чл.~31 ЗЕДЕУУ).
Надзорът върху доставчиците упражнява КРС (чл.~32, ал.~1 ЗЕДЕУУ).

За лицето, на чието име е издадено удостоверение въз основа на чужда лична карта,
презумпцията за вина на квалифицирания доставчик е най-силната защита. В делото, в което
е представен документът, то носи тежестта да докаже, че не е подписало. В отделен иск
срещу доставчика обаче е достатъчно да докаже вредата и неизпълнението на задължението за
проверка на самоличността, а доставчикът трябва да докаже, че не е виновен. Така рискът е
разпределен между три лица: титулярят носи риска от доброволното предоставяне на подписа и
от неуведомяването за промени, адресатът~-- риска от липсата на дължима грижа, а
доставчикът~-- риска от неправилното издаване на удостоверението.
